\documentclass[12pt,a4paper,openright]{report}

\usepackage[T1]{fontenc}
\usepackage[utf8]{inputenc}
\usepackage[english]{babel}
\usepackage[a4paper,margin=1in]{geometry}
\usepackage{setspace}
\usepackage{graphicx}
\usepackage{booktabs}
\usepackage{array}
\usepackage{longtable}
\usepackage{amsmath,amssymb}
\usepackage{hyperref}
\usepackage{float}
\usepackage[numbers,sort&compress]{natbib}

\hypersetup{colorlinks=true, linkcolor=black, citecolor=black, urlcolor=blue}
\onehalfspacing

\title{Fact Validator: Transparent, Source-Aware Misinformation Detection and Fact Verification}
\author{Candidate Name}
\date{Academic Year 2025--2026}

\begin{document}

\begin{titlepage}
\begin{center}
{\Large UNIVERSITA DEGLI STUDI DI NAPOLI FEDERICO II\par}
\vspace{0.5cm}
{\large Department of Electrical Engineering and Information Technologies\par}
\vspace{0.3cm}
{\large Master's Degree Programme in Computer Engineering / Artificial Intelligence\par}
\vspace{1.7cm}
{\Large \textbf{Master's Thesis}\par}
\vspace{1.2cm}
{\LARGE \textbf{Fact Validator: Transparent, Source-Aware Misinformation Detection and Fact Verification}\par}
\vspace{1.2cm}
{\large A systems-and-methodology thesis on open-domain fact checking, credibility-aware evidence retrieval, and deployment-ready misinformation analysis\par}
\vfill
\begin{minipage}[t]{0.45\textwidth}
\raggedright
\textbf{Supervisor}\\
Prof. Supervisor Name\\[0.4cm]
\textbf{Co-supervisor}\\
Dr. Co-supervisor Name
\end{minipage}
\hfill
\begin{minipage}[t]{0.45\textwidth}
\raggedleft
\textbf{Candidate}\\
Candidate Name\\
Matricola: 000000
\end{minipage}
\vfill
{\large Academic Year 2025--2026\par}
\end{center}
\end{titlepage}

\pagenumbering{roman}

\chapter*{Abstract}
This thesis presents Fact Validator, a deployment-oriented system for fake-news and misinformation analysis. The repository contains two related but separate artifacts. The Fact Validator live application accepts URLs or raw text, retrieves evidence, applies auditable domain credibility scoring, semantically reranks evidence, and persists user-facing reports. The frozen 5000-claim experiment evaluates a deterministic FactValidator-Proxy that combines lexical classification, category-level priors, heuristic semantic signals, deterministic arbitration, and a quality filter. It does not execute SerpAPI, live domain scoring, SentenceTransformer reranking, or Ollama debate for every test claim. On FEVER, LIAR, SciFact, and the PUBHEALTH \texttt{health\_fact} dataset, the full proxy reaches 50.82\% accuracy and 0.4384 macro-F1; the no-debate proxy reaches 51.34\% and 0.4427. The novelty claim concerns transparent, auditable, deployment-aware systems integration rather than universal superiority in factuality classification.

\chapter*{Thesis Claim}
The core thesis claim is that a fact-checking system can be made more publishable and practically useful when it treats source credibility, evidence persistence, benchmarking reproducibility, and operational cost as first-class design dimensions. Fact Validator demonstrates this claim through an auditable architecture and a versioned evaluation workflow.

\tableofcontents
\clearpage

\pagenumbering{arabic}

\chapter{Introduction}

Misinformation detection is not only a classification problem. A realistic system must decide what to verify, where to search for evidence, how much to trust each source, how to explain a verdict, and how to stay usable when external tools fail. Most benchmark-centered papers optimize one slice of this workflow. This thesis studies the live architecture required for practical fact checking and evaluates its separate deterministic proxy without treating the two evidence surfaces as identical.

The goal of the live Fact Validator application is to assess the likelihood that content is fake or misleading by combining retrieval, credibility reasoning, claim-level verification, and a user-facing analysis interface. Its implementation evidence is kept separate from the deterministic proxy benchmark used for the 5000-claim statistical evaluation. The motivating research question is the following:

\begin{quote}
Can an end-to-end fact-checking architecture achieve a meaningful balance between factuality quality, explainability, and deployment realism without hiding source trust inside a black-box model?
\end{quote}

\section{Research Objectives}

This thesis pursues five objectives:

\begin{enumerate}
  \item design a full-stack pipeline for claim extraction and evidence-backed fact verification,
  \item make source trust explicit through an auditable credibility prior,
  \item compare the resulting system to relevant literature in faithful correction and tool-augmented factuality,
  \item measure proxy predictive behavior and report separately labelled operational projections,
  \item produce a publication-ready research narrative with honest limitations.
\end{enumerate}

\chapter{State of the Art}

\section{Zero-shot Faithful Factual Error Correction}

Zero-shot Faithful Factual Error Correction addresses factual error detection and rewriting in a zero-shot regime. Its strongest idea is that correction quality should remain faithful to the source content and grounded in external knowledge. This is important for the present thesis because it shows that factuality can be handled without narrow task-specific retraining. However, its scope is narrower than misinformation analysis over articles or multi-claim source text. Fact Validator inherits the emphasis on faithfulness, but extends the workflow to evidence retrieval, domain trust scoring, and final misinformation assessment.

\section{FacTool}

FacTool is the closest conceptual reference because it formalizes factuality detection as a tool-augmented process. The key overlap is the use of external retrieval and modular verification components. The main difference is that Fact Validator places more emphasis on auditable source credibility, persistent storage, fallback logic, and cost accounting. FacTool is a framework for tool-enabled factuality analysis; Fact Validator is a research system intended to operate as a practical misinformation-checking platform.

\section{Thesis Positioning}

The contribution of this thesis is strongest when framed along three axes:

\begin{enumerate}
  \item \textbf{Scope:} from sentence-level correction or factuality checking to complete misinformation assessment.
  \item \textbf{Transparency:} explicit source-credibility scoring rather than hidden trust estimation.
  \item \textbf{Deployment realism:} caching, persistence, fallback behavior, and reproducible evaluation artifacts.
\end{enumerate}

\chapter{System Architecture}

\section{High-Level Pipeline}

Fact Validator uses the following nine-stage architecture:

\begin{enumerate}
  \item content extraction from a URL or raw text input,
  \item claim decomposition into fact-like propositions,
  \item web evidence retrieval,
  \item domain credibility scoring,
  \item semantic reranking,
  \item verdict classification into SUPPORTED, REFUTED, or NEI,
  \item optional debate arbitration,
  \item sentiment and bias analysis,
  \item final misinformation likelihood scoring and persistence.
\end{enumerate}

\section{Layered View}

\begin{table}[H]
\centering
\caption{Layered architecture of Fact Validator.}
\begin{tabular}{p{3cm}p{9.5cm}}
\toprule
Layer & Main responsibility \\
\midrule
Frontend & Next.js interface for URL or text submission, result review, and research dashboard access \\
API orchestration & FastAPI service coordinating extraction, scoring, verification, export, and persistence \\
Verification core & claim decomposition, retrieval, credibility scoring, semantic reranking, verdict inference, optional debate \\
Storage & SQLite run history, JSON caches, reproducibility artifacts \\
External services & SerpAPI, optional Ollama endpoint, transformer-based similarity models \\
\bottomrule
\end{tabular}
\end{table}

\section{Credibility Scoring as a Thesis Contribution}

The credibility prior is a central novelty claim of this thesis. Instead of learning source trust implicitly, Fact Validator exposes a score that can be inspected and disputed:

\[
\mathrm{Credibility}(d) = \mathrm{clip}_{[0,100]}\bigl(50 + b(d) + p(d)\bigr)
\]

where $b(d)$ provides bonuses for reliable domains and $p(d)$ applies penalties for low-trust domains. This makes the system more suitable for academic defense because the trust model is falsifiable.

\section{Why This Architecture Is Strong}

The architecture is strong in ways that matter for a thesis committee:

\begin{enumerate}
  \item it is end-to-end rather than a disconnected benchmark script,
  \item it integrates explanation and source trust into inference,
  \item it stores evidence and outputs for later audit,
  \item it measures cost and latency explicitly,
  \item it remains usable without the optional debate layer.
\end{enumerate}

\chapter{Methodology and Experimental Design}

\section{Benchmark Selection}

The repository contains several result families. For a defensible final thesis narrative, this document uses the strongest frozen benchmark family: the 5000-claim evaluation generated on 2026-07-02. That benchmark combines FEVER, LIAR, SciFact, and a normalized health-domain input source. Using one frozen run family avoids the confusion caused by smaller pilot splits and older draft results.

\section{Reported Variants}

The evaluation includes the full deterministic proxy, tuned and ablated proxy variants, and simple baselines. This allows the thesis to separate three questions:

\begin{enumerate}
  \item How does the full proxy compare with naive comparators?
  \item Which internal components matter?
  \item Which parts help operational realism without necessarily increasing accuracy?
\end{enumerate}

\section{Metrics}

The thesis reports:

\begin{enumerate}
  \item accuracy and macro-F1,
  \item raw heuristic score and the correctly scaled raw-score diagnostic,
  \item per-dataset deltas,
  \item latency, throughput, and monthly cost estimates.
\end{enumerate}

\chapter{Results}

\section{Main Accuracy Results}

\begin{table}[H]
\centering
\caption{Main results on the frozen 5000-claim benchmark.}
\begin{tabular}{lrrrr}
\toprule
System & Accuracy & 95\% Wilson CI & Avg. raw score & Raw-score gap \\
\midrule
ablate\_debate & 0.5134 & [0.4995, 0.5273] & 88.37 & 0.3703 \\
tune\_fever & 0.5098 & [0.4959, 0.5237] & 88.53 & 0.3755 \\
full\_proxy & 0.5082 & [0.4943, 0.5221] & 88.31 & 0.3749 \\
length & 0.4942 & [0.4803, 0.5081] & 48.83 & 0.0059 \\
majority & 0.4926 & [0.4787, 0.5065] & 50.00 & 0.0074 \\
random & 0.3320 & [0.3189, 0.3451] & 50.00 & 0.1680 \\
\bottomrule
\end{tabular}
\end{table}

The full proxy exceeds the majority baseline by 1.56 percentage points and the length heuristic by 1.40 percentage points. Exact unadjusted McNemar p-values are 0.0433 and 0.0462, so these marginal comparisons are not described as robust superiority after multiple-comparison correction. The single highest frozen result is the no-debate proxy at 51.34\%; this finding concerns deterministic proxy arbitration, not live Ollama debate.

\section{Per-Dataset Behavior}

\begin{table}[H]
\centering
\caption{Per-dataset performance of the full deterministic proxy.}
\begin{tabular}{lrrrrr}
\toprule
Dataset & $n$ & Full & Majority & Length & Full-Majority \\
\midrule
FEVER & 1829 & 0.5752 & 0.5856 & 0.5790 & -0.0104 \\
Health domain & 1490 & 0.5517 & 0.5154 & 0.5255 & +0.0362 \\
LIAR & 1490 & 0.3859 & 0.3658 & 0.3698 & +0.0201 \\
SciFact & 191 & 0.4817 & 0.4136 & 0.4084 & +0.0681 \\
\bottomrule
\end{tabular}
\end{table}

This table shows that proxy behavior varies by dataset: its largest descriptive margins occur on health and scientific claims, while it is slightly below simple baselines on FEVER. Because the proxy does not execute live retrieval and domain scoring for every item, these results do not establish that source-sensitive retrieval caused the differences.

\section{Operational Results}

\begin{table}[H]
\centering
\caption{Operational scenario projections; not controlled load-test measurements.}
\begin{tabular}{lr}
\toprule
Metric & Value \\
\midrule
Baseline latency assumption & 8.20 s/claim \\
Debate latency assumption & 72.00 s/claim \\
Projected debate overhead & 8.78x \\
Projected baseline throughput & 439.02 claims/hour \\
Projected debate throughput & 50.00 claims/hour \\
Projected monthly cost without caching & USD 77.00 \\
Projected monthly cost with caching & USD 22.00 \\
Projected monthly savings & 71.43\% \\
\bottomrule
\end{tabular}
\end{table}

These scenarios are part of the research value because they make cost assumptions explicit. They are not a substitute for repeated warm- and cold-cache load tests.

\chapter{Discussion}

\section{Why the Research Is Better, and in What Sense}

The thesis should not claim that Fact Validator is universally better than Zero-shot Faithful Factual Error Correction or FacTool in overall factuality quality. That would exceed the evidence. The stronger and defensible statement is narrower:

\begin{enumerate}
  \item it places greater emphasis on a deployment-ready misinformation-analysis surface,
  \item it makes source-trust assignment explicit and inspectable,
  \item it documents operational assumptions and reproducibility artifacts,
  \item it contributes a broad end-to-end research artifact.
\end{enumerate}

\section{Novelty}

The novelty of this thesis lies in the combination of:

\begin{enumerate}
  \item auditable credibility scoring,
  \item claim-level verification over live web evidence,
  \item optional live debate in the application and separately evaluated deterministic arbitration in the proxy,
  \item persistent evidence, storage, and export surfaces,
  \item versioned large-scale evaluation with operational reporting.
\end{enumerate}

This is a stronger novelty claim than simply saying the model is more accurate.

\section{Limitations}

Three limitations must be stated clearly in the thesis defense:

\begin{enumerate}
  \item the proxy's raw confidence score remains much higher than its observed accuracy and is not presented as a calibrated probability,
  \item debate does not improve the frozen benchmark when always enabled,
  \item the health-domain slot uses a normalized substitute source and should be documented transparently.
\end{enumerate}

\chapter{Conclusion}

Fact Validator demonstrates that open-domain fact checking can be framed as a transparent systems problem rather than only a benchmark optimization problem. The project contributes a full architecture for evidence-backed misinformation analysis, a clear source-credibility prior, a reproducible evaluation pipeline, and an honest comparative narrative against strong reference papers. The best supported final thesis claim is therefore this: Fact Validator advances the state of practical, auditable fact verification by integrating verification quality, explainability, and operational realism in one reproducible system.

\appendix

\chapter{Defense-Ready Summary Paragraph}

Fact Validator should be defended as a transparent and deployment-ready misinformation analysis architecture with a separate deterministic proxy evaluation. The proxy shows modest raw differences over majority and length baselines, but neither comparison survives Holm correction. Its larger scientific value lies in explicit source-trust modeling, persistent evidence workflows, reproducible artifacts, and clearly labelled cost scenarios. Compared with Zero-shot Faithful Factual Error Correction and FacTool, the system is most novel not because it dominates their native tasks, but because it operationalizes factuality research as a practical, auditable, end-to-end platform.

\begin{thebibliography}{99}

\bibitem{kubota2023zero}
Kubota, S., Kajiwara, Y., and Onishi, T. (2023).
Zero-shot Faithful Factual Error Correction.
In \textit{Proceedings of ACL 2023}.

\bibitem{factool2023}
FacTool: Factuality Detection in Generative AI -- A Tool-Augmented Framework.
2023.
\textit{Conference publication}.

\bibitem{thorne2018fever}
Thorne, J., Vlachos, A., Christodoulopoulos, C., and Mittal, A. (2018).
FEVER: A Large-scale Dataset for Fact Extraction and Verification.
In \textit{Proceedings of NAACL-HLT}, pages 809--819.

\bibitem{ji2023hallucination}
Ji, Z., Lee, N., Frieske, R., et al. (2023).
Survey of Hallucination in Natural Language Generation.
\textit{ACM Computing Surveys}, 55(12), 1--38.

\bibitem{wen2024abstention}
Wen, B., Yao, J., Feng, S., et al. (2024).
Know Your Limits: A Survey of Abstention in Large Language Models.
\textit{arXiv preprint arXiv:2407.18418}.

\end{thebibliography}

\end{document}
